\chapter{Considerações Finais} \label{cpt:discus}
Este capítulo resume nossas contribuições para o desenvolvimento de pesquisa no uso de tecnologias para a área da saúde, as limitações da abordagem proposta e nossos planos futuros para continuidade da pesquisa.

\section{Contribuições}
Neste trabalho desenvolvemos um conjunto de dados com informações reais de quantidade de leitos hospitalares separados por especialidade e tipo nos hospitais de Porto Alegre, em complemento também adicionamos informações geográficas de latitude e longitude de cada hospital. Os dados coletados foram provenientes do Cadastro Nacional de Estabelecimentos de Saúde (CNES) disponibilizado pelo DATASUS e as informações geográficas coletadas no Google Maps. O conjunto de dados gerado pode ser utilizado em outras aplicações voltadas para a área da saúde ou gerenciamento urbano ou hospitalar. 

Além do conjunto de dados, também foi proposta uma ferramenta de simulação de ocupação de leitos hospitalares na cidade de Porto Alegre chamado de LODUS-Health. Foi disponibilizado uma interface gráfica para facilitar a interação dos usuários com o modelo de simulação, permitindo parametrizações customizadas e integração dos dados do conjunto de dados criado inicialmente, além da possibilidade de geração de dados médicos para uma determinada especialidade de leito a partir da seleção do usuário na interface. Analisando os resultados, conclui-se que ferramenta desenvolvida pode ser utilizada por gestores de hospitais do município na destinação de recursos para cada hospital ou ate mesmo simular a inclusão de um novo hospital em um determinado bairro.

Visto que o LODUS-Health é um módulo do LODUS, cada vez que um paciente chega ao hospital, ele saiu de casa, se transportou de alguma maneira e foi ao hospital. Isso dá outra perspectiva à simulação urbana que permite fechar-se, abrir-se novos hospitais, e as pessoas se dirigirem aos estabelecimentos.

\section{Limitações}
O modelo de simulação desenvolvido e apresentado nesse trabalho está limitado a um contexto onde espera-se que o usuário da ferramenta possua conhecimento prévio sobre o fluxo de ocupação de leitos hospitalares para que consiga configurar de forma mais realista os parâmetros da simulação. %Além disso, em casos de superlotação dos leitos hospitalares, não há uma diferenciação entre os pacientes que estão nos leitos e os que estão esperando por vagas, apenas sabe se da quantidade de pacientes requisitados para aquele leito. 
Outra limitação da pesquisa é a necessidade da coleta de dados sobre os leitos dos hospitais precisarem ser feitas de forma recorrente devido a atualização dessas informações no CNESNet. Isso é relevante para manter os dados atualizados com a cidade que se deseja simular. Apesar disso, é importante mencionar que a quantidade de leitos não muda bruscamente de uma atualização para a outra, logo atualizações são necessárias, mas a hipótese é de que não sejam tão frequentes. %, logo a liié necessário essa coleta para que a simulação tenha os dados o mais próximo da realidade atual.

\section{Trabalhos Futuros}
A abordagem utilizada no desenvolvimento gerou uma gama de possíveis trabalhos futuros como complemento dessa pesquisa e desenvolvimento da área. Algumas sugestões para trabalhos futuros incluem: \textit{i)} Prover uma visualização relevante das movimentações dos pacientes na cidade e nos hospitais, \textit{ii)} Permitir que o gestor escolha dentre os pacientes, aquele que se deseja ver o texto médico gerado,  \textit{iii)} Aprimorar o controle da fila do hospital para permitir troca de leitos, pessoas que são transferidas para outro hospital, \textit{iv)} associar um CID especifico a um paciente para geração de texto dessa doença e etc.


% coletar dados reais de internacao

% permitir a troca de pacientes entre hospitais

% permitir interacao entre outros tipos de pontos de interesse

% pacientes individuais